\subsection{Engineering Problem}

Reliable autonomous flight requires more than recognizing that a multivariate telemetry sequence has departed from normal operation. A binary anomaly decision is valuable for surveillance and safety monitoring, and real-time methods have demonstrated that abnormal UAV behavior can be detected from onboard flight data \cite{keipour2019automatic,wei2026aerotsboost}. In a flight-control engineering workflow, however, the same alert must also support a decision about what to inspect. The engineer needs to know which fault domain is implicated, which telemetry channels provide the strongest evidence, and which parts of the flight-control software consume or produce the corresponding signals. This distinction reflects the broader separation between fault detection and fault diagnosis in UAV systems \cite{fourlas2021survey,puchalski2022uav,tan2026anomaly}.

The distinction is especially important for heterogeneous PX4 logs. Position estimates, estimator innovations, actuator responses, inertial measurements, battery states, and system-health variables are coupled through the closed-loop vehicle dynamics. Consequently, a high anomaly score alone cannot distinguish an external-position disturbance from a global-position, altitude, or mechanical/electrical anomaly. These domains imply different follow-up checks even when their telemetry deviations overlap. The UAV-SEAD corpus provides real PX4 flight logs with these runtime anomaly categories \cite{kabaoglu2026uavsead}; our data audit retained 1,389 usable logs, establishing a sufficiently broad but class-imbalanced basis for evaluating both binary detection and subsequent fault-domain diagnosis.

Actionable diagnosis also requires an interpretable path from model output to the flight-control architecture. PX4 is organized as a set of modules that exchange typed messages through its publish--subscribe middleware \cite{meier2015px4,px4uorbdocs2026}. A telemetry channel can therefore be associated with a logged uORB topic and a functional subsystem, while a publisher/subscriber graph reconstructed for a specified firmware revision can identify software modules connected to that topic. Such a chain turns statistical evidence into architecture-aware inspection leads. It does not, by itself, prove that any connected module caused the observed behavior; the resulting modules are suspects to prioritize during engineering inspection rather than localized source-code faults.

This evidentiary boundary is central to the problem formulation. Labels attached to real flights describe observed runtime conditions or fault domains, not defective PX4 statements, functions, or modules. Treating those labels as source-code ground truth would collapse two different inference tasks: diagnosing vehicle behavior from telemetry and localizing an injected or known software defect. We therefore formulate the engineering problem as a hierarchy: first detect abnormal telemetry, then diagnose the affected fault domain, expose the telemetry evidence supporting that decision, and finally use PX4 architecture only to rank plausible module inspection leads. Evidence for the first two stages is provided by the audited real-flight corpus and the detection and fault-domain experiments, whereas claims about software-module ranking require separate controlled evidence.
